% TOP_PROBLEM: {"id": 7500109, "title": "Reverse-Mathematical Strength of Hindman's Theorem", "category_id": 18, "category": {"id": 18, "name": "logic", "display_name": "Logic", "description": "Problems in mathematical logic, model theory, proof theory, and finite model theory.", "slug": "logic", "order_index": 18, "created_at": "2026-07-31T15:26:25.671Z"}, "status": "partially_solved"}
% FIRST_POSTED: null
% ENTRY_KIND: statement_only
% Imported from: batch20.tex; UnsolvedMath ID 7500109
\documentclass[11pt]{article}
\usepackage[a4paper,margin=25mm]{geometry}
\usepackage{amsmath,amssymb,mathtools}
\usepackage{fontspec,unicode-math}
\setmainfont{Latin Modern Roman}
\setsansfont{Latin Modern Sans}
\setmathfont{Latin Modern Math}
\usepackage{xurl,hyperref}
\hypersetup{colorlinks=true,linkcolor=blue,urlcolor=blue,pdftitle={UnsolvedMath Problem 7500109}}
\newcommand{\sref}[2]{\hyperlink{source:#1:#2}{[S#2]}}
\newcommand{\cataloguescope}[1]{\paragraph{Mathematical question.}#1}
\begin{document}
% UnsolvedMath ID 7500109; source catalogue code AMR-074-0109
\setcounter{section}{7500108}
\section{Reverse-Mathematical Strength of Hindman's Theorem}\label{7500109}
{\small\sffamily UnsolvedMath ID 7500109 · Logic}
\subsection{Definitions and mathematical statement}

\paragraph{Shared notation.}$\mathbb N=\{1,2,\ldots\}$, and $\mathbb Z,\mathbb Q,\mathbb R,\mathbb C$ denote the integers, rationals, reals and complex numbers. Any other conventions are specified locally.


Work in second-order arithmetic over $\mathsf{RCA}_0$, the base theory with $\Delta^0_1$ comprehension and $\Sigma^0_1$ induction. The strength of a principle means which subsystems prove it and which subsystems it implies over that base; equivalence requires both implications.  Here $\mathsf{ACA}_0$ adds comprehension for all arithmetical formulas. For a set $X$, its Turing jump $X'$ codes which oracle computations with oracle $X$ halt; $X^{(\omega)}=\bigoplus_{n\geq0}X^{(n)}$ joins all finite jumps. The theory $\mathsf{ACA}_0^+$ is $\mathsf{RCA}_0+\forall X\,\exists X^{(\omega)}$. Hindman’s theorem $\mathsf{HT}$ says that, for every finite colouring of $\mathbb N$, some infinite $A\subseteq\mathbb N$ has all nonempty finite sums of distinct members of $A$ of one colour.
\paragraph{Statement recorded in the supplied source.}
Is $\mathsf{HT}$ equivalent to $\mathsf{ACA}_0^+$, equivalent to $\mathsf{ACA}_0$, or strictly between these two theories? \sref{7500109}{2}
\subsection{Short English statement}
Does Hindman’s finite-sums theorem require exactly arithmetic comprehension, exactly the existence of all finite jumps joined together, or a strength strictly between them?
\subsection{Sources}
\begin{description}
\item[\hypertarget{source:7500109:1}{S1}] ULAM.AI and the UnsolvedMath Contributors, UnsolvedMath: A Curated Collection of Open Mathematics Problems, record 7500109 (AMR-074-0109). CC BY 4.0. \url{https://huggingface.co/datasets/ulamai/UnsolvedMath}.
\item[\hypertarget{source:7500109:2}{S2}] Antonio Montalbán, Open Questions in Reverse Mathematics, author preprint dated 17 August 2010, Question 9, printed p. 5. \url{https://math.berkeley.edu/~antonio/papers/questionsRM.pdf}.
\end{description}
\label{end:7500109}
\end{document}
